% [Mahmoud] Results and Discussion (rubric: 3 marks), with paragraphs from Sumaiya (leakage audit,
% stage analysis) and Lynn (promise test).
% Top band: clear labelled figures and tables; FAIR comparison against the baseline; error analysis and
%   interpretation of what the model learned; robustness; whether differences are MEANINGFUL (bootstrap CIs).
% Rubric note: a report whose only result is a single model score cannot reach the top band.
% Say whether each proposal expectation held (distance/peaks/load top drivers; -20% MAE; errors up at peaks).
% All numbers from results/final/ (CP4), never v0.
\section{Results and Discussion}

\subsection{Model comparison}
TODO. % Table: test MAE/RMSE/R^2/late rate for median, promise, Elastic-Net, RF, boosting + bootstrap CIs
% Final test table (CP4), generated by experiments.final; refer to it as Table~\ref{tab:main_results}.
% The same table on the CV folds is tables/main_results_cv (label tab:main_results_cv) if needed.
\begin{table*}[t]
\centering
\caption{Delivery-time prediction on the test days (23--24 Oct). $\Delta$MAE in minutes with 95\% paired bootstrap intervals; negative = lower error. Late $>$5 min: arrival more than 5 minutes after the promise (platform) or after the model's estimate (other rows).}
\label{tab:main_results}
\begin{tabular}{lrrrrll}
\toprule
Model & MAE & RMSE & $R^2$ & Late $>$5 min & $\Delta$MAE vs median & $\Delta$MAE vs promise \\
\midrule
Platform promise & 9.62 & 11.87 & 0.048 & 5.0\% & +0.16 [+0.11, +0.21] & -- \\
Median predictor & 9.46 & 12.38 & -0.035 & 36.1\% & -- & -0.16 [-0.21, -0.11] \\
Elastic-Net & 6.80 & 8.98 & 0.455 & 27.6\% & -2.66 [-2.69, -2.62] & -2.82 [-2.86, -2.77] \\
Random Forest & 6.54 & 8.67 & 0.492 & 26.9\% & -2.92 [-2.95, -2.88] & -3.08 [-3.12, -3.04] \\
Gradient boosting & 6.49 & 8.70 & 0.488 & 29.4\% & -2.97 [-3.00, -2.93] & -3.13 [-3.17, -3.08] \\
\bottomrule
\end{tabular}
\end{table*}


\subsection{Delay drivers}
TODO. % grouped permutation importance (M5)

\subsection{Where the model fails}
TODO. % errors by district, hour, distance, test day (M6)

\subsection{Leakage audit and delivery stages}
TODO. % [Sumaiya] S5, S6

\subsection{Promise setting}
% [Lynn] L4: q80 vs platform at the same mean promised time
A point estimate aims at the middle of the distribution, so about half of all orders arrive after it;
a promise needs a higher quantile. Our 80th-percentile model, shifted by $+2.7$ minutes so that its
average promise equals the platform's (38.2 minutes on the test days), covers 84.7\% of orders against
the platform's 85.1\%. At this equal average, the platform is late slightly less often at every
threshold (Table~\ref{tab:promise}): 15.3\% vs 14.9\% of orders arrive after the promise, 6.1\% vs 5.0\%
more than five minutes late, and 2.4\% vs 1.5\% more than ten minutes late, all with bootstrap intervals
that exclude zero. Fig.~\ref{fig:promise_gap} shows that the comparison depends on promise length:
with shorter average promises (below about 37 minutes) our model is late less often at the zero-minute
threshold, but at the platform's average and above, the platform does better. On the validation folds our model was marginally better at
the zero-minute threshold ($-0.3$ points) but already worse beyond five minutes, so the pattern is
stable: the platform is better at preventing \emph{very} late orders. A likely reason is that it can
still act after the order is placed. The EDA points the same way: the number of orders arriving late
more than halves at exactly eight minutes late, at every promise level, which looks like an operational
rule applied after ordering (Section~\ref{sec:data}). A model
limited to order-time information cannot do this, so better estimates (Table~\ref{tab:main_results})
do not by themselves give better promises.
% Table~\ref{tab:promise} and Fig.~\ref{fig:promise_gap}; CV version: tables/promise_test_cv
\begin{table}[t]
\centering
\caption{Promise setting at the same mean promised time (38.2 min; our q80 promises shifted by +2.7 min). Test days. Gap = ours $-$ platform, 95\% bootstrap CI.}
\label{tab:promise}
\begin{tabular}{lrrl}
\toprule
Metric & Ours & Platform & Gap \\
\midrule
Late $>$0 min & 15.3\% & 14.9\% & +0.3 pp [+0.2, +0.5] \\
Late $>$5 min & 6.1\% & 5.0\% & +1.2 pp [+1.1, +1.3] \\
Late $>$10 min & 2.4\% & 1.5\% & +0.9 pp [+0.9, +1.0] \\
Pinball loss ($q=0.8$) & 2.52 & 2.36 & -- \\
Coverage & 84.7\% & 85.1\% & -- \\
\bottomrule
\end{tabular}
\end{table}

\begin{figure}[t]
\centering
\includegraphics[width=\columnwidth]{final_test_gap}
\caption{Late-rate gap (our q80 promises minus the platform's, percentage points) as both are shifted to the same mean promised time; below 0 means we are late less often. Dotted line: the platform's actual mean promise. Test days.}
\label{fig:promise_gap}
\end{figure}
